\subsection{通过平展且满射的映射下降}

设 X 和 Y 是拓扑空间，f 是从 X 到 Y 的连续映射。沿用前一节的记号。

定理 2. --- 假设 Y 的每一点都有一个邻域，在其上存在映射 f 的连续截面。那么，群胚态射 \(\varpi'(f)\) 是从 \(\operatorname{Coeg}(f)\) 到 \(\varpi(Y)\) 的同构。

根据假设，存在由开集构成的 Y 的覆盖 \((\mathrm{U}_{j})_{j\in\mathrm{J}}\)，并且对每个 \(j\in J\)，有连续截面 \(s_{j}\)，它是 \(f_{U_{j}}\) 的截面。

若 \(c\) 是 X 中的一条道路，我们将 \([c]\) 记为它在 \(\varpi(X)\) 中的严格同伦类，将 \(\{c\}\) 记为 \([c]\) 在 \(\operatorname{Coeg}(f)\) 中经群胚态射 \(\gamma\) 得到的像。若 \(c\) 和 \(c'\) 是 X 中的两条道路，则 \(\{c\}\) 和 \(\{c'\}\) 在 \(\operatorname{Coeg}(f)\) 中可合成，当且仅当道路 \(f \circ c\) 和 \(f \circ c'\) 在 Y 中可首尾相接。

设 \(c'\) 是 Y 中的一条道路。根据 III, 第~272~页引理 4，将其应用于紧空间 I 以及 I 的覆盖 \(((c')^{-1}(\mathrm{U}_j))_{j\in \mathrm{J}}\)，存在整数 \(n\)，使得对每个整数 \(k\)，若满足 \(1\leqslant k\leqslant n\)，则区间 \([\frac{k - 1}{n},\frac{k}{n}]\) 在 \(c'\) 下的像包含于某个开集 \(\mathrm{U}_{j(k)}\) 中。对每个整数 \(k\)，若 \(1\leqslant k\leqslant n\)，令 \(c_k'\) 为 Y 中由 \(s\mapsto c'(\frac{k + s - 1}{n})\) 定义的道路；有 \([c'] = [c_1'][c_2']\ldots [c_n']\)（参见 III, 第~291~页，注 1），且 \(c'\) 就是记作 \(c_1' * c_2' * \dots * c_n'\) 的道路。对所有 \(k\in \{1,\ldots ,n\}\)，记 \(c_k\) 为 X 中的道路 \(s_{j(k)}\circ c_k'\)，并置 \(\{c_k\} = \gamma ([c_k])\)。由于对每个 \(k\)，道路 \(c_{k - 1}' = f\circ c_{k - 1}\) 和 \(c_k' = f\circ c_k\) 可首尾相接，序列 \((\{c_1\} ,\ldots ,\{c_n\})\) 在 \(\operatorname{Coeg}(f)\) 中可合成。由构造，

\[
\begin{array}{r l} & {\varpi^ {\prime} (f) (\{c _ {1} \} \ldots \{c _ {n} \}) = \varpi^ {\prime} (f) (\{c _ {1} \}) \ldots \varpi^ {\prime} (f) (\{c _ {n} \})} \\ & {\qquad = \varpi (f) ([ c _ {1} ]) \ldots \varpi (f) ([ c _ {n} ])} \\ & {\qquad = [ f \circ c _ {1} ] \ldots [ f \circ c _ {n} ]} \\ & {\qquad = [ c _ {1} ^ {\prime} ] * \dots * [ c _ {n} ^ {\prime} ] = [ c ^ {\prime} ],} \end{array}
\]

这证明群胚态射 \(\varpi'(f)\) 是满射。

设 \(u\) 和 \(v\) 是 \(\operatorname{Coeg}(f)\) 的箭。由于群胚 \(\operatorname{Coeg}(f)\) 由 \(\varpi(\mathrm{X})\) 的像生成（II, 第~200~页，推论），存在 X 中道路的有限序列 \((c_1, \ldots, c_n)\) 和 \((d_1, \ldots, d_n)\)，使得 \(u = \{c_1\} \ldots \{c_n\}\) 且 \(v = \{d_1\} \ldots \{d_n\}\)。道路 \((f \circ c_1, \ldots, f \circ c_n)\) 在 Y 中可首尾相接，并且有 \(\varpi'(f)(u) = [(f \circ c_1)] \ldots [(f \circ c_n)]\)；同样，\(\varpi'(f)(v) = [(f \circ d_1)] \ldots [(f \circ d_n)]\)。

假设有 \(\varpi'(f)(u) = \varpi'(f)(v)\)。于是存在严格同伦 \(\sigma\)，连接 \((f \circ c_1) * \cdots * (f \circ c_n)\) 与 \((f \circ d_1) * \cdots * (f \circ d_n)\)。根据 III, 第~272~页引理 4，将其应用于紧空间 \(\mathbf{I} \times \mathbf{I}\) 以及覆盖 \((\bar{\sigma}^{-1}(\mathrm{U}_i))_{i \in \mathrm{I}}\) 的空间 \(\mathbf{I} \times \mathbf{I}\)，存在整数 \(m \geqslant 1\)，使得对每一对整数 \((j,k)\)，若满足 \(1 \leqslant j \leqslant m\) 且 \(1 \leqslant k \leqslant m\)，则区间积 \([\frac{j-1}{m}, \frac{j}{m}] \times [\frac{k-1}{m}, \frac{k}{m}]\) 在 \(\sigma\) 下的像包含于某个开集 \(U_{i(j,k)}\) 中，该开集属于覆盖 \((\mathrm{U}_{i})_{i \in \mathrm{I}}\)。

X 中每条道路 c 都形如 \(c_{1}*\cdots*c_{m}\)，其中 \(c_{k}\) 是由 \(t\mapsto c(\frac{k-1+t}{m})\) 定义的道路。如有必要，将整数 m 和 n 换成它们的乘积 mn，因此可以假设 m=n。

对于每一对整数 \((j,k)\)（属于 \(\{1,\ldots ,n\}\)）以及每一对 \((s,t)\in\) \(\mathbf{I}\times \mathbf{I}\)，置

\[
\sigma_ {j, k} (s, t) = s _ {i (j, k)} \circ \sigma \left(\frac {s + j - 1}{n}, \frac {t + k - 1}{n}\right).
\]

对于 \(t \in \mathbf{I}\)，再置 \(h_{j,k}^{0}(t) = \sigma_{j,k}(t,0)\)、\(h_{j,k}^{1}(t) = \sigma_{j,k}(t,1)\)、\(v_{j,k}^{0}(t) = \sigma_{j,k}(0,t)\) 以及 \(v_{j,k}^{1}(t) = \sigma_{j,k}(1,t)\)。根据 III, 第~295~页引理 1，道路 \(h_{j,k}^{0} * v_{j,k}^{1}\) 和 \(v_{j,k}^{0} * h_{j,k}^{1}\) 严格同伦，因此有关系

\[
[ h _ {j, k} ^ {0} ] [ v _ {j, k} ^ {1} ] = [ v _ {j, k} ^ {0} ] [ h _ {j, k} ^ {1} ]\tag{2}
\]

在 \(\varpi(\mathrm{X})\) 上，对每一对 \((j,k) \in \{1, \ldots, n\}^2\) 都成立。另一方面，对每一对整数 \((j,k)\)，若 \(2 \leqslant j \leqslant n\) 且 \(1 \leqslant k \leqslant n\)，有 \(f \circ v_{j,k}^0 = f \circ v_{j-1,k}^1\)，因此有关系

\[
\{v _ {j, k} ^ {0} \} = \{v _ {j - 1, k} ^ {1} \}\tag{3}
\]

在 Coeg(f) 中。类似地，对每一对整数 \((j,k)\)，若 \(1 \leqslant j \leqslant n\) 且 \(2 \leqslant k \leqslant n\)，有

\[
\{h _ {j, k} ^ {0} \} = \{h _ {j, k - 1} ^ {1} \}.\tag{4}
\]

对于 \(j \in \{1, \ldots, n\}\)，道路 \(f \circ c_j\) 和 \(f \circ h_{j,1}^0\) 重合。根据余均衡子 Coeg(f) 的定义，因此有 \(\{c_j\} = \{h_{j,1}^0\}\)。同样，对每个 \(j \in \{1, \ldots, n\}\)，有 \(\{d_j\} = \{h_{j,n}^1\}\)。因此有如下关系

\[
u = \{h _ {1, 1} ^ {0} \} \dots \{h _ {n, 1} ^ {0} \}, \quad v = \{h _ {1, n} ^ {1} \} \dots \{h _ {n, n} ^ {1} \}.
\]

根据下面的引理 3，有

\[
\{h _ {1, 1} ^ {0} \} \dots \{h _ {1, n} ^ {0} \} \{v _ {1, n} ^ {1} \} \dots \{v _ {n, n} ^ {1} \} = \{v _ {1, 1} ^ {0} \} \dots \{v _ {n, 1} ^ {0} \} \{h _ {n, 1} ^ {1} \} \dots \{h _ {n, n} ^ {1} \}.
\]

由于道路 \(t \mapsto \sigma(0,t)\) 和 \(t \mapsto \sigma(1,t)\) 是常值，对 \(1 \leqslant k \leqslant n\)，有 \(\{v_{1,k}^{0}\} = e_{a}\) 且 \(\{v_{n,k}^{1}\} = e_{b}\)，其中 \(a,b \in Y\) 是箭 \(u\) 和 \(v\) 的起点和终点，它们属于群胚 Coeg( \(f\) )。因此有等式

\[
\{h _ {1, 1} ^ {0} \} \dots \{h _ {1, n} ^ {0} \} = \{h _ {n, 1} ^ {0} \} \dots \{h _ {n, n} ^ {1} \},
\]

即 u = v。

引理 3. --- 设 G 是群胚，p、q 是满足 \(\geqslant 1\) 的整数。对于每一对整数 \((j, k)\)，若 \(0 \leqslant j \leqslant p\) 且 \(0 \leqslant k \leqslant q\)，令 \(x_{j,k}\) 为 G 的一个顶点；对于每一对整数 \((j, k)\)，若 \(1 \leqslant j \leqslant p\) 且 \(0 \leqslant k \leqslant q\)，令 \(h_{j,k}\) 为连接 \(x_{j-1,k}\) 与 \(x_{j,k}\) 的 G 的一个箭；对于每一对整数 \((j, k)\)，若 \(0 \leqslant j \leqslant p\) 且 \(1 \leqslant k \leqslant q\)，令 \(v_{j,k}\) 为连接 \(x_{j,k-1}\) 与 \(x_{j,k}\) 的 G 的一个箭。

假设对于每一对整数 \((j,k)\)，若 \(1 \leqslant j \leqslant p\) 且 \(1 \leqslant k \leqslant q\)，箭 \(v_{j-1,k}\) 和 \(h_{j,k}\) 可合成，箭 \(h_{j,k-1}\) 和 \(v_{j,k}\) 也可合成，并且有 \(v_{j-1,k}h_{j,k} = h_{j,k-1}v_{j,k}\)。那么，

\[
h _ {1, 0} h _ {2, 0} \dots h _ {p, 0} v _ {p, 1} v _ {p, 2} \dots v _ {p, q} = v _ {0, 1} v _ {0, 2} \dots v _ {0, q} h _ {1, q} h _ {2, q} \dots h _ {p, q}.
\]

先处理 q = 1 的特殊情形，并对 p 作归纳证明。若 p = 1，待证断言由假设得到验证；假设它对 p - 1 成立；于是有

\[
h _ {1, 0} h _ {2, 0} \dots h _ {p, 0} v _ {p, 1} = h _ {1, 0} h _ {2, 0} \dots h _ {p - 1, 0} v _ {p - 1, 1} h _ {p, 1} = v _ {0, 1} h _ {1, 1} \dots h _ {p, 1}
\]

根据归纳假设，从而得到 p 的关系。

现在对 q 作归纳证明。根据前述内容，q = 1 时命题成立；若对 q - 1 成立，则有

\[
h _ {1, 0} h _ {2, 0} \dots h _ {p, 0} v _ {p, 1} v _ {p, 2} \dots v _ {p, q} = v _ {0, 1} v _ {0, 2} \dots v _ {0, q - 1} h _ {1, q - 1} \dots h _ {p, q - 1} v _ {p, q}.
\]

根据 \(q = 1\) 的情形，有

\[
h _ {1, q - 1} \dots h _ {p, q - 1} v _ {p, q} = v _ {0, q} h _ {1, q} \dots h _ {p, q},
\]

从而得到所需关系。

例。--- 1) 当映射 f 平展且满射时，定理适用。

\begin{enumerate}
\def\labelenumi{\arabic{enumi})}
\setcounter{enumi}{1}
\tightlist
\item
  它也适用于如下情形：空间 X 是集合族 \((V_{i})_{i\in\mathrm{I}}\) 的和空间，其中这些 Y 的子集的内部覆盖 Y，且 f 是由每个 \(V_{i}\) 到 Y 的规范嵌入导出的映射。
\end{enumerate}

